\documentclass[11pt]{article}
\usepackage[T1]{fontenc}
\usepackage[utf8]{inputenc}
\usepackage{mathptmx}
\usepackage[margin=1in]{geometry}
\usepackage{amsmath,amssymb}
\usepackage[hidelinks]{hyperref}
\usepackage{enumitem}
\usepackage{microtype}
\usepackage{tabularx}
\usepackage{booktabs}
\PassOptionsToPackage{hyphens}{url}
\usepackage{url}
\makeatletter
\g@addto@macro\UrlBreaks{\do\/\do\_\do\.\do\-}
\makeatother
\Urlmuskip=0mu plus 1mu\relax
\usepackage[dvipsnames]{xcolor}
\definecolor{ReviewerGreen}{HTML}{008000}
\newcommand{\rvcomment}[1]{{\color{ReviewerGreen}#1}}
\newcommand{\concern}[1]{\textbf{\color{ReviewerGreen}#1}}
\setlength{\parindent}{0pt}
\setlength{\parskip}{6pt}
\setlist[itemize,1]{leftmargin=1.6em,labelsep=0.6em,itemsep=3pt,topsep=3pt,parsep=0pt,
                    label=\raisebox{-0.15ex}{\large\textbullet}}
\setlist[itemize,2]{leftmargin=1.4em,labelsep=0.5em,itemsep=2pt,topsep=3pt,parsep=0pt,
                    label=\textendash}
\pagestyle{plain}
\begin{document}

\concern{Reviewer \#2, Concern \#1}

\concern{Reviewer comment:}

\rvcomment{1. Please formalize the ``near-boundary'' component of the operational review rule in
Section 4 and Algorithm 1. The function $g$ maps a letter rating to a decision class, such as
IG/HY, and $\hat r$ is subsequently defined as the modal decision class. It is therefore unclear how
the condition $\hat r \in B$ identifies a rating ``near'' the IG/HY boundary.}

\textbf{Author response:}

\begin{itemize}
\item We agree that the original rule did not formally represent proximity to the IG/HY boundary. Because
$g(r)$ maps a letter rating to a decision class, the modal class $\hat r$ can only be IG or HY and
therefore cannot encode distance from the boundary. We have corrected the rule so that boundary
proximity is defined using the modal letter rating and its notch distance from the BBB$-$/BB$+$
boundary.
\end{itemize}

\textbf{Author action:}

\begin{itemize}
\item Algorithm 1 now defines boundary proximity from the modal letter rating using the half-step boundary $b=\{\mathrm{rank}_L(\mathrm{BBB}-)+\mathrm{rank}_L(\mathrm{BB}+)\}/2$. The operational referral rule is stated separately on the dissent count, $d\ge c_{\mathrm{ref}}$. Because instability is not confined to the boundary region, proximity is returned as a priority indicator rather than a mandatory referral condition. Requiring both dissent and $w=1$ proximity would refer 13.3\% of profiles with recall 0.176.
\item The earlier ``roughly halves to about 33\%'' statement has been removed; on the held-out split, referring on any disagreement ($c_{\mathrm{ref}}=1$) covers 88.9\% of profiles.
\item We have removed the earlier statement that $w=1$ reproduces the WATCH stratum. WATCH is the Altman $Z''$ stratum spanning BB$+$ through B$-$, a different construct, whereas the rule measures distance from the BBB$-$/BB$+$ cut. The manuscript now calls this an operational near-boundary band, and the stale Section~3.5 cross-reference has been deleted.
\end{itemize}

\vspace{4pt}

\concern{Reviewer \#2, Concern \#2}

\concern{Reviewer comment:}

\rvcomment{2. Please clarify the sampling convention underlying the reported matched-pair flip
probabilities. The Methods and the captions of Figures 2--4 should state whether two draws may
select the same elicitation cell and whether the pair may differ in specification, seed, or both.
This clarification would make the 25.2\% result and the self-consistency analysis fully
reproducible.}

\textbf{Author response:}

\begin{itemize}
\item We agree that the sampling convention should have been stated explicitly. The reported flip
probability,
\[
1-\sum_c p_c^2,
\]
is the probability that two independent draws with replacement from the same firm's observed
elicitation distribution fall in different classes. The firm profile is always held fixed. The two draws may differ in
specification, in seed, or in both. Identical-cell draws are permitted and count as agreement.
\item A complete credit-health profile contributes $32\times3=96$ elicitations, though four profiles are short one seed and contribute $95$, so we no longer write $n=96$ universally. Using draws without replacement, evaluated at each profile's own $n_i$ rather than with a blanket $96/95$, changes the pooled estimate only from $25.2\%$ to 25.45\%.
\item We also label the $25.2\%$ explicitly as total elicitation instability, since the two draws may differ in specification, in seed, or in both, and report two complementary conditional measures, which are not additive: 24.3\% after collapsing repeated seeds within specification and 6.9\% when varying seed within an identical profile$\times$specification cell. The four two-observation cells are unanimous, so the modal tie rule does not bind. The headline figure is not attributed to specification choices alone.
\end{itemize}

\textbf{Author action:}

\begin{itemize}
\item We added the full sampling convention to Section 6.3 and to the caption of Figure 2, covering replacement, same-cell treatment, the permitted sources of variation, and the fixed-firm condition.
\item We removed it from Figure 3, which reports Fleiss $\kappa$ and is not computed from sampled pairs, and Figure 4's caption now states the procedure actually used: $p$ is the profile's empirical investment-grade share, the $k$ votes are modelled as independent Bernoulli$(p)$ draws so $q$ follows in closed form from the binomial, ties at even $k$ are resolved to either class with probability one half, and the plotted residual is the mean of $2q(1-q)$.
\item Section 3.5 now reports the overall 25.2\% elicitation-disagreement rate together with these two complementary conditional measures.
\end{itemize}

\vspace{4pt}

\concern{Reviewer \#2, Concern \#3}

\concern{Reviewer comment:}

\rvcomment{3. Please harmonize the description of the analytical sample near the beginning of the
paper. The abstract reports 45 credit-health profiles, 32 specifications, and three seeds, giving
4,320 credit-rating elicitations, whereas the Introduction describes a 90-item battery and 8,640
elicitations. Although Section 6.1 explains that the additional 45 directional items are not
analyzed in this paper, this distinction should be stated at the first mention of the 8,640-cell
experiment to avoid the impression of inconsistent sample sizes.}

\textbf{Author response:}

\begin{itemize}
\item We agree that the analytical and full-battery sample counts should be reconciled at their first
appearance. The full battery contains two families of 45 profiles. This paper analyses the 45
credit-health profiles, corresponding to $45\times32\times3=4{,}320$ designed elicitation cells, of
which $4{,}316$ returned substantive content. The additional 45 directional profiles use the same
specification grid, are included in the archive, and are not analysed in this paper.
\end{itemize}

\textbf{Author action:}

\begin{itemize}
\item The Introduction now makes this distinction at the first mention of the full experiment, and the manuscript consistently distinguishes designed from realized sample sizes.
\item The Methods note that four of the seven non-substantive cells occur in the credit-health arm. These cells were not re-collected, because the corpus had already been frozen.
\end{itemize}

\vspace{4pt}

\concern{Reviewer \#2, Concern \#4}

\concern{Reviewer comment:}

\rvcomment{4. In Table 4 and Section 3.6, the statement that approximately 33\% agreement is
``expected by chance'' requires qualification. If this is a uniform three-class reference value,
please label it as such. If it is intended as a chance-corrected baseline based on the observed
class marginals, the corresponding calculation should be provided.}

\textbf{Author response:}

\begin{itemize}
\item We agree that the uniform $33.3\%$ reference is not the appropriate chance benchmark for the
observed class distributions. Using the empirical machine and agency marginals, expected agreement
is $48.9\%$, compared with observed agreement of $53.7\%$. The excess agreement is therefore $4.9$
percentage points (issuer-clustered $95\%$ CI $-7.9$ to $17.1$), with Cohen's $\kappa=0.095$.
\item Accordingly, the real-issuer arm does not provide evidence that machine--agency agreement exceeds
the marginal-matching chance benchmark. This correction does not affect the separate within-issuer
specification contrast reported for this arm.
\end{itemize}

\textbf{Author action:}

\begin{itemize}
\item Section 3.6 and Table 4 now report observed agreement, marginal-matching expected agreement, the excess with its confidence interval, and Cohen's $\kappa$.
\item The manuscript no longer interprets machine--agency agreement as evidence of rating accuracy.
\end{itemize}

\vspace{4pt}

\concern{Reviewer \#3, Concern \#1}

\concern{Reviewer comment:}

\rvcomment{the authors have made efforts to address the previous review comments, including
reframing the contribution as an evaluation-methodology case study, removing the economic
translations, correcting the permutation test, adding a reasoning-architecture arm, and providing
held-out-specification validation for the SIS\dots\ however, the revised manuscript still requires
careful evaluation, and the following points must be considered before it can be further considered
for publication. this manuscript is a research article, and it is evaluated as such.}

\rvcomment{the per-axis variance decomposition remains non-identified under the
resolution-limited d-optimal fractional grid, yet the abstract, highlights, and conclusions still
state that no controllable axis exceeds about 5\% of rating variance. the authors now report
aliasing diagnostics, including a maximum absolute correlation of 0.333 between a main effect and a
two-factor interaction of other factors, and they correctly label the decomposition as descriptive.
however, the headline claim of ``no axis exceeds \textasciitilde5\%'' is still presented as a
substantive finding rather than as a model-dependent descriptive partition. because main effects can
be partially aliased with omitted interactions, the per-axis point estimates cannot support even a
descriptive ranking of axes without explicit bounds or a design that separates them. the added
R\textsuperscript{2} range of 0.083 to 0.119 across model specifications does not identify individual
axes; it only shows that the elicitation axes jointly explain little variance under the assumed
models. this remains important because the paper's central interpretation---that instability is
spread across elicitation choices rather than concentrated in a tunable knob---depends in part on
this decomposition.}

\rvcomment{required action: either remove all per-axis numeric claims from the abstract, highlights,
and conclusions, or conduct a fully crossed or resolution-v sub-experiment on a reduced set of
factors and items that can cleanly estimate main effects and relevant two-way interactions. if a new
experiment is not feasible, provide formal bounds on each per-axis variance share under all aliasing
structures consistent with the observed design matrix, or fit a bayesian hierarchical model with
explicit shrinkage priors on interaction terms and report posterior intervals for the main effects.
the abstract and highlights should then state only what is identified: that the elicitation axes
jointly explain a small and model-dependent share of variance, not that any specific axis is bounded
by \textasciitilde5\%.}

\textbf{Author response:}

\begin{itemize}
\item We agree that the fractional design does not identify individual elicitation-axis variance shares
once two-way interactions are allowed. We have therefore removed the per-axis interpretation rather
than attempting to rank individual axes. Formally, if $Z$ denotes the omitted interaction contrasts,
\[
\mathbb{E}[\hat{\beta}]=\beta+A\gamma,\qquad A=(X'X)^{-1}X'Z,
\]
so that $|a_j'\gamma|\leq\lVert a_j\rVert_2\lVert\gamma\rVert_2$. Because the design imposes no
finite restriction on $\lVert\gamma\rVert_2$, no finite design-only bound exists for an individual
main effect. For answer presentation $\lVert a_j\rVert_2=0.878$, so the former ${\sim}4.6\%$
estimate cannot support an identified per-axis interpretation.
\item We also recomputed the aliasing diagnostics on the 32-specification design actually used in the
variance decomposition. The maximum main-effect--interaction correlation is $0.504$, and
canonical-correlation analysis shows substantial overlap between the main-effect and interaction
subspaces, with the largest canonical correlation at $1.000$ and six of eight above $0.9$. As a
quality-control check, we re-audited all design-dependent diagnostics and variance-decomposition
inputs on the analysed 32-specification arm, and no other design-dependent diagnostic or variance-decomposition result changed. The
paper therefore no longer makes any claim about how instability is distributed across individual
elicitation axes. The retained result is only that, under the fitted models, the elicitation axes
jointly account for $R^2=0.083$--$0.119$ of observed variation. The model-free disagreement
results are unaffected by this decomposition.
\item The provider$\times$item term is different because items are crossed with specifications. All
$45\times32=1{,}440$ item--specification cells are represented, with $1{,}436$ containing all three
seeds and four containing two. The resulting imbalance is numerically negligible, because the
largest realized correlation with a specification-axis contrast is only $0.0022$. We therefore describe the
provider$\times$item component as identified up to this small imbalance.
\end{itemize}

\textbf{Author action:}

\begin{itemize}
\item We removed the ``no individual axis exceeds ${\sim}5\%$'' statement from the Abstract, Highlights, Results and Discussion.
\item Section 6.2 now states explicitly what the design does and does not identify, including the aliasing algebra, corrected design diagnostics and canonical-correlation analysis.
\item Individual-axis estimates have been removed from Table 1.
\end{itemize}

\vspace{4pt}

\concern{Reviewer \#3, Concern \#2}

\concern{Reviewer comment:}

\rvcomment{the real-issuer arm still does not provide strong external validity, and the
decontamination procedure remains insufficiently validated for the claims made\dots\ the perturbation
applies a single per-issuer uniform scaling factor to all dollar quantities. this preserves financial
ratios and Altman Z'' by construction, but it does not prevent an LLM from inferring issuer identity
from the pattern of ratios, nor does it rule out memorized issuer knowledge that is not based on
exact figures. the fingerprinting test only measures exact-figure lookup. the authors acknowledge
this limitation, but the abstract still states that specification variation is observed in ``31
perturbed and fingerprint-screened real issuers,'' which can be read as external validation. in
addition, the scale factor is confounded with issuer identity because each issuer receives only one
draw. the sample of 31 issuers and 12 specifications is also small for any quantitative external
claim.}

\rvcomment{required action: add a ratio-pattern identification test. for example, train a classifier
on financial ratios to predict issuer identity from the perturbed profiles, or prompt the LLM to
infer the issuer from the ratio profile after perturbation, and report the identification rate.
conduct a within-issuer scale sensitivity design by re-eliciting a subset of issuers at multiple
scale factors, such as 0.6, 1.0, and 1.7, so that scale is not fully confounded with issuer
identity. if these checks show residual identification or if the sample remains too small, remove
the real-issuer claim from the abstract and conclusion and present the arm only as an illustrative
descriptive limitation. the paper should not imply external validity beyond what the design can
support.}

\textbf{Author response:}

\begin{itemize}
\item We agree that uniform scaling alone cannot establish decontamination, because it preserves
financial ratios and the original design assigned only one scale factor per issuer. We therefore
performed both additional checks suggested by the reviewer, under a specification fixed and hashed
before the new model calls.
\item First, we conducted a within-issuer scale experiment. All 31 issuers were re-elicited at scales
$0.6$, $1.0$ and $1.7$ on the frozen 12-specification grid, making each issuer its own control. The
IG/HY class differed between the extreme scales in $15.9\%$ of issuer--specification cells,
compared with $15.2\%$ under the within-issuer permutation null ($p=0.644$, 10{,}000 permutations).
We therefore do not detect evidence that absolute scale contributes additional instability beyond
the re-elicitation variability represented by this null. The experiment uses one realization per
issuer--specification--scale cell and is therefore not intended to estimate seed-level variance
separately. Its purpose is narrower, namely to remove the original confounding between issuer
identity and scale and to test scale labels within matched issuer--specification blocks.
\item Second, we conducted a ratio-pattern identification experiment. In the open-ended probe, neither
provider correctly identified an issuer in the 60 substantive completions from 62 attempted calls.
In the more stringent closed-set probe, however, the correct issuer was selected in $19.4\%$ of
substantive calls ($95\%$ CI $12.4\%$--$27.4\%$), substantially above the $1/31=3.2\%$ chance
benchmark ($p=5\times10^{-18}$), and nineteen of the 31 issuers were identified correctly at least
once. Because the ratio-identification check detects residual identifying
information, we adopt the conservative branch proposed by the reviewer. The 31-issuer arm is
retained only as an illustrative descriptive check and is not used as evidence of external
validity. The closed-set probe should itself be interpreted conservatively, because supplying the
candidate names creates a stronger identification setting than an open-ended task. Taken together,
the new checks resolve the design question without upgrading the arm's evidentiary status. We find
no detectable incremental scale effect under the paired test, but ratio patterns retain enough
identifying information that ratio-preserving scaling cannot be treated as decontamination.
\end{itemize}

\textbf{Author action:}

\begin{itemize}
\item We added Section 3.7 and Table 5 reporting both new experiments.
\item The Abstract, Section 3.6 and Conclusion now describe the real-issuer arm as illustrative and descriptive, and make no external-validity claim from it.
\item The pre-registered experimental specification, frozen outputs, manifests and analysis code have been added to the archive.
\end{itemize}

\vspace{4pt}

\concern{Reviewer \#3, Concern \#3}

\concern{Reviewer comment:}

\rvcomment{the SIS validation is improved, but its incremental value over simple
disagreement-based baselines and its status as a deployable model-risk control remain unproven\dots\
at the binary IG/HY gate with a low threshold, the SIS flag reduces to ``any disagreement'' across
the k elicitations. the reported Youden-optimal threshold of 0.08 is below the smallest non-zero SIS
value at k=16, so the operating rule is essentially ``flag any issuer with at least one disagreeing
calibration elicitation.'' this rule is already directly observable from the k elicitations and does
not require the continuous SIS. the manuscript does not compare SIS against simple baselines such as
the number of dissenting elicitations, the modal rating's distance to the IG/HY boundary,
self-consistency vote margins, semantic entropy, or token-level confidence. the threshold is also
chosen and evaluated on the same held-out split, and cross-model transfer is moderate and
inconsistent. this matters because the abstract and highlights present the SIS as a ``validated'' and
``deployable model-risk control.''}

\rvcomment{required action: benchmark SIS against these simple baselines in both
held-out-specification and held-out-issuer settings. report a decision-curve analysis or net-benefit
analysis for triage, showing whether SIS improves referral decisions over ``flag any disagreement,''
over distance to the IG/HY boundary, and over self-consistency vote margins. if SIS does not provide
incremental value beyond these baselines, reframe it as a descriptive dispersion index rather than a
deployable control, and remove or substantially qualify the ``validated'' and ``deployable'' language
in the abstract and highlights. if it does provide incremental value, report the threshold selection
in a fully nested cross-validation or on an independent calibration set and show how performance
varies with k and tau.}

\textbf{Author response:}

\begin{itemize}
\item We agree that the relevant question is whether $\mathrm{SIS}$ provides information beyond simpler
measures of disagreement, and we therefore benchmarked it directly. As requested, we evaluated
these comparisons in both the held-out-specification setting and a profile-split subsampling analysis, reported
separately in Table 8. The same conclusion holds in both, namely that $\mathrm{SIS}$ provides no
incremental ranking information over the corresponding dissent count at the binary IG/HY gate. In
a profile-split subsampling analysis, which splits the 45 credit-health profiles at random into
halves over 300 draws and scores only the held-out half, the mean AUC is 0.936 at the gate and 0.950 at three bands, and the
index-minus-dissent-count difference is exactly zero on every split. At the binary IG/HY
resolution,
if $m$ of $k$ elicitations belong to the minority class,
\[
\mathrm{SIS}=2(m/k)(1-m/k),
\]
which is strictly increasing in $m$ for $m<k/2$. $\mathrm{SIS}$, dissent count and vote margin are
therefore rank-equivalent at the binary gate, and the empirical results confirm this exactly.
Each yields AUC $0.936$, the bootstrap AUC difference is $0.000$, the Spearman correlation is
$1.000$, and their net-benefit curves coincide. At three bands the measures are not rank-identical, with a Spearman correlation of $0.971$ between the index and the dissent count, yet they yield the same AUC in the present specification-split and profile-split analyses. The operational rule is therefore stated directly on the dissent count, $d\ge c_{\mathrm{ref}}$, with the threshold chosen according to referral cost and capacity. On the held-out split, $c_{\mathrm{ref}}=3$ yields precision 1.000 and recall 0.882 and dominates the $c_{\mathrm{ref}}=1$ ``any disagreement'' rule when review is sufficiently costly. The former Youden-selected $\tau=0.08$ is retained nowhere in the operating recommendation; for the record it equals $d\ge3$ at the held-out $k=48$ and would collapse to $c_{\mathrm{ref}}=1$ at $k=16$.
\item In preparing this response, we identified and corrected a numerical inconsistency. Section 4.2 previously reported precision $1.00$ and recall $0.88$ for ``any disagreement,'' which cannot hold alongside a net benefit that varies with the threshold probability. Both are now generated from one calculation. Referring on any disagreement gives $\mathrm{TP}=33$, $\mathrm{FP}=7$, precision $0.825$ and recall $0.971$, whereas the $1.00$/$0.88$ pair belongs to $\tau=0.08$ on the same split, which gives $\mathrm{TP}=30$, $\mathrm{FP}=0$ and hence a flat net benefit of $0.6667$ at every threshold probability. The two rules had been conflated in the prose, not miscomputed.
\item We also benchmarked boundary distance and discrete semantic entropy. Boundary distance is
materially weaker. For discrete semantic entropy, following Farquhar et al.\ (2024), probabilities
can be estimated from empirical generation frequencies without token probabilities. Because our
outputs are categorical ratings, the corresponding discrete entropy reduces to the Shannon entropy
of the rating distribution, which is again rank-equivalent at the binary gate and produces the same
AUC and net benefit as $\mathrm{SIS}$ and dissent count. Two scope limits apply. Farquhar et al.\
(2024) study repeated generations under a fixed prompt, whereas our application varies elicitation
specifications, so we use this only as a specification-level analogue of the estimator. Token-level
confidence cannot be reconstructed, because the frozen corpus retains no log-probabilities and
these are not uniformly available across the providers evaluated. Because $\mathrm{SIS}$ shows no incremental value and the manuscript no longer proposes a tuned
$\mathrm{SIS}$ threshold as a predictive instrument, we do not estimate or validate an optimized
$\tau$. Nested threshold selection would answer a deployment question that the revised manuscript
no longer claims to answer. The resulting interpretation is simpler, in that early specification
disagreement predicts later disagreement, but $\mathrm{SIS}$ does not provide incremental
predictive information beyond the corresponding dissent count at the binary gate. We therefore present $\mathrm{SIS}$ as a
descriptive dispersion index and governance representation rather than as an incrementally
validated predictive instrument.
\end{itemize}

\textbf{Author action:}

\begin{itemize}
\item Section 4.3 and Table 8 now benchmark $\mathrm{SIS}$ against dissent count, vote margin, boundary distance and discrete semantic entropy, in both the held-out-specification setting and a profile-split subsampling analysis, reported separately, and a separate net-benefit analysis following Vickers and Elkin (2006) compares the corresponding referral strategies with refer-all and refer-none in Table 9.
\item The Abstract and Highlights no longer describe $\mathrm{SIS}$ as a validated or deployable model-risk control.
\item Sections 4.1--4.3 instead present it as a descriptive dispersion index and governance trigger.
\end{itemize}

\vspace{4pt}

\concern{Reviewer \#3, Concern \#4}

\concern{Reviewer comment:}

\rvcomment{under the resolution-limited d-optimal fractional grid, can you provide formal
identifiable bounds on the maximum possible main effect of answer presentation and the
provider-by-item interaction under all aliasing structures consistent with the observed design
matrix, and how would those bounds change the abstract's claim that no controllable axis exceeds
about 5\% of rating variance?}

\textbf{Author response:}

\begin{itemize}
\item The two quantities in the reviewer's question have different identification properties. For answer
presentation, no finite design-only bound exists. With omitted interaction coefficients $\gamma$,
\[
\mathbb{E}[\hat{\beta}]=\beta+A\gamma,
\]
and the alias contribution is bounded only conditional on an external restriction on
$\lVert\gamma\rVert_2$, which the observed design does not provide. The former ${\sim}4.6\%$
answer-presentation estimate therefore cannot be interpreted as an identified upper bound or used
to rank that axis against others.
\item For provider$\times$item, the analysed $45\times32$ item--specification grid is fully represented,
although four cells contain two rather than three seeds. The resulting imbalance is very small, with the largest realized
correlation between provider$\times$item and a specification-axis contrast at only $0.0022$. We therefore describe the approximately $1.4\%$ provider$\times$item component as
identified up to this numerically negligible imbalance.
\item The implication for the headline result is that the former statement that no controllable axis
exceeds ${\sim}5\%$ has been removed. The manuscript now reports only the joint explanatory share
under the fitted specifications ($R^2=0.083$--$0.119$), without assigning that share to individual
elicitation axes. Thus the answer to the reviewer's bounds question changes the abstract
qualitatively rather than numerically. There is no defensible finite upper bound supporting the
former ${\sim}5\%$ statement, so that statement has been removed rather than weakened.
\end{itemize}

\textbf{Author action:}

\begin{itemize}
\item Section 6.2 now provides the formal aliasing analysis and numerical diagnostics.
\item The Abstract and Highlights no longer contain the per-axis ${\sim}5\%$ statement, and individual-axis estimates have been removed from Table 1.
\end{itemize}

\vspace{4pt}

\end{document}
